\chapter{Introduction}

SAT solvers, in my opinion, are one of the most underrated pieces of algorithmic technology of the modern day. They are used in both software and hardware verification (circuitry testing), product configuration, package management, computational biology, cryptanalysis, particle physics, solving graph problems, and much much more.\footnote{\url{https://www.cs.utexas.edu/~isil/cs389L/lecture4-6up.pdf}}

In a more computer science-focused world, they have applications in NP-Completeness. Since (nearly) every problem has a reduction to a Boolean SAT problem, creating one really really good SAT solver lets us solve all those problems in a ``fast'' manner. So they are important everywhere despite the vast majority of us never really thinking about their existence.

One note on what ``fast'' means: even if solving a SAT problem takes a small amount of time, the algorithmic complexity isn't necessarily ``fast'' as there is no known polynomial algorithm for any NP-Complete problem (I am a firm believer that $P = NP$, and yes I know that makes me crazy).

Before we start writing out code, proofs, and the like, we need to make sure we're all on the same page on things such as input and ideas. Reading through this introduction is important so that you understand the basic notation used throughout this book.

One last thing to note: we will write out pseudocode rather than the code for one specific programming language. However, although the idea behind pseudocode is that it's programming language agnostic (it doesn't carry artifacts of any real programming language (usually)), this idea is poorly executed when it comes to formatting.

My algorithms professor will probably scream at me for this, but I'm making the executive decision that we will follow a general C-based syntax for readability purposes. That is, we will use braces \{\} for block scope, parentheses () for grouping, \verb|=| for assignment, \verb|==| for equality, and \verb|<=|, \verb|>=|, etc.\ for their respective uses. Everything that can be inferred visually will be used.

Another thing to note is that we will pretend our arrays can dynamically resize, so we won't pre-allocate space, and we will also assume 1-based indexing for an easier map from variables to indices, as you will see later.

\section{What Do SAT Solvers Even Solve?}

``What do SAT solvers even solve,'' I hear you say? SAT solvers solve Boolean satisfiability problems. Another way to say that is they solve propositional logic problems. There are many forms of propositional logic, and you may be familiar with it. Here are some basic examples:

\begin{itemize}
    \item $P \rightarrow Q$, an implication statement
    \item $(A \land B)$, an and statement
    \item $(A \lor B)$, an or statement
    \item $\neg A$, a negation/not statement
\end{itemize}

Of course these are very basic examples, and we haven't included every operator such as xor, nand, etc. As you know, we can rewrite implications to contain only and, or, and not operators, which is what SAT solvers operate on. In other words, SAT solvers only operate on Boolean algebra equations.\footnote{A Boolean algebra equation is one containing only and ($\land$), or ($\lor$), and not ($\neg$) operations (at a basic level).}

However, SAT solvers don't operate on just any Boolean algebra equation.

No no no, they must be in a proper form, called CNF (more on that in the next section). CNF (Conjunctive Normal Form) equations are a ``conjunction of disjunctions,'' i.e., clauses of ors, combined with ands.

Some examples of CNF are as follows:

\begin{itemize}
    \item $(A \lor B)$
    \item $(A \lor B) \land (\neg C)$
    \item $(A \lor B) \land (\neg C \lor D)$
\end{itemize}

Some examples that are \textbf{not} CNF:

\begin{itemize}
    \item $\neg(A \lor B)$, where there is a NOT in front of a clause
    \item $\neg(A \lor B) \land (C)$, where there is a NOT in front of a clause
    \item $(A \lor B \land C)$, where there is an AND within a clause
    \item $(A \land B) \lor (C)$, where it is in DNF form.\footnote{DNF (Disjunctive Normal Form) is a ``disjunction of conjunctions,'' i.e., clauses of ands combined by ors. Solving these is trivially easy, as we just need to find one clause where no element evaluates to false.}
\end{itemize}

\section{CNF Encoding}

Now that we know what our inputs look like, the next question (I hope you're
asking questions) is ``how do we encode that?'' Well good thing you're reading
this, because that's what we're here to talk about! In this section, we'll expand
our base of boolean algebra, including some definitions, transformations, and
finally good encoding.

\subsection*{Definitions We'll Use}

\begin{description}
    \item[Clause] A clause is a collection of variables, in which each variable is
    separated by an or ($\lor$), and each clause is separated by an and
    ($\land$).

    \item[Variable] A variable is an element of a clause, taking on a value of
    True, False, or Unassigned. Each variable can be negated ($\neg$) or not,
    affecting its evaluated value.

    \item[Literal] A literal is an evaluated variable, evaluating to one of True,
    False, or Unassigned. Since a literal is an evaluated variable, it cannot be
    negated. A literal whose corresponding variable is unassigned evaluates to
    Unassigned, regardless of whether or not the variable is negated.
\end{description}

A quick example of the difference of a literal value, using all three of our
definitions:

\[
(\neg A)
\]

is a clause of one variable ``A'', which is negated. If the variable
$A=\text{False}$, then its literal value is $\text{True}$, because

\[
\neg \text{False} \equiv \text{True}.
\]

In other words, $A$ was assigned the value False, and then it evaluated to
True. This will be very important in the future, so please take the time to
understand it well enough.

As discussed earlier, a CNF equation only has three operators: and ($\land$),
or ($\lor$), and not ($\neg$), but in propositional logic, there are many more
operators, and their clauses are not guaranteed to be in the format we want.
There are many algorithms to transform any propositional logic equation into
CNF, but frankly, that's beyond the scope of this book (it's also beyond my
scope). Instead, we will just look at how we transform each operator into an
equivalent format using just our three operators.

\subsection*{Implication ($\rightarrow$)}

\begin{center}
\begin{tabular}{cc|c|c|c}
$p$ & $q$ & $\neg p$ & $p\rightarrow q$ & $\neg p\lor q$\\
\hline
T & T & F & T & T\\
T & F & F & F & F\\
F & T & T & T & T\\
F & F & T & T & T
\end{tabular}
\end{center}

\subsection*{Xor ($\oplus$)}

\begin{center}
\begin{tabular}{cc|cc|c|cc|c}
$p$ & $q$
& $\neg p$
& $\neg q$
& $p\oplus q$
& $p\lor q$
& $\neg p\lor\neg q$
& $(p\lor q)\land(\neg p\lor\neg q)$\\
\hline
T&T&F&F&F&T&F&F\\
T&F&F&T&T&T&T&T\\
F&T&T&F&T&T&T&T\\
F&F&T&T&F&F&T&F
\end{tabular}
\end{center}

\subsection*{Biconditional ($\leftrightarrow$)}

\begin{center}
\begin{tabular}{cc|c|cc|c}
$p$ & $q$
& $p\leftrightarrow q$
& $p\rightarrow q$
& $q\rightarrow p$
& $(\neg p\lor q)\land(\neg q\lor p)$\\
\hline
T&T&T&T&T&T\\
T&F&F&F&T&F\\
F&T&F&T&F&F\\
F&F&T&T&T&T
\end{tabular}
\end{center}

We won't look at negations of operators such as nand, nor, xnor, or any
others. I'm not a boolean algebra expert when it comes to what operators
exist. But if you were given an equation with these extra operators, you could
easily apply the correct equivalency, then apply an algorithm that converts any
boolean algebra equation into one of CNF form.\footnote{I honestly may be
talking out of my ass here. I will do more research on this subject at a later
date.}

\subsection*{DIMACS CNF Format}

Now that we know what CNF looks like, we can now discuss how it's formatted
for a SAT solver's consumption. Typically, CNF is formatted in
\emph{DIMACS CNF} form.

Some rules for DIMACS CNF form are:

\begin{enumerate}
    \item The file may begin with comment lines. The first character of each
    comment line must be a lower case letter ``c''. Comment lines typically
    occur in one section at the beginning of the file, but are allowed to appear
    throughout the file.

    \item The comment lines are followed by the ``problem'' line. This begins
    with a lower case ``p'' followed by a space, followed by the problem type,
    which for CNF files is ``cnf'', followed by the number of variables followed
    by the number of clauses.

    \item The remainder of the file contains lines defining the clauses, one by
    one.

    \item A clause is defined by listing the index of each positive literal, and
    the negative index of each negative literal. Indices are 1-based, and for
    obvious reasons the index 0 is not allowed.

    \item The definition of a clause may extend beyond a single line of text.

    \item The definition of a clause is terminated by a final value of
    ``0''.

    \item The file terminates after the last clause is defined.
\end{enumerate}

There are some things about this syntax to be wary of:

\begin{enumerate}
    \item The definition of the next clause normally begins on a new line, but
    may follow, on the same line, the ``0'' that marks the end of the previous
    clause.

    \item In some examples of CNF files, the definition of the last clause is
    not terminated by a final ``0''.

    \item In some examples of CNF files, the rule that the variables are
    numbered from 1 to $N$ is not followed. The file might declare that there
    are 10 variables, for instance, but allow them to be numbered 2 through
    11.

    \item Comments can appear anywhere in the file, but the comment must be
    on its own line and it must start with ``c''.\footnote{\url{https://jix.github.io/varisat/manual/0.2.0/formats/dimacs.html}}
\end{enumerate}

For a quick example, here is a simple way to format a DIMACS CNF equation:

\begin{lstlisting}
p cnf 4 2
1 2 -3 0
-1 4 0
\end{lstlisting}

And one with comments:

\begin{lstlisting}
c this is a comment
p cnf 4 2
c this is another comment
1 2 -3 0
-1 4 0
\end{lstlisting}

In order to increase the compatibility of our SAT solver, we will be bending
some rules.

We will firstly allow missing header lines. In the case where no header is
passed in, we will instead infer the number of clauses and variables. We will
also be very lenient in where comments can appear, but we will enforce that
variables must be between $1$ and $N$.

\begin{enumerate}
    \item Comments can appear anywhere in the file, but they \textbf{MUST}
    be on their own line and start with ``c''.

    \item A header is not necessary, and if one is not provided, then the
    number of variables and clauses shall be inferred. However, if one is
    provided, then the number of variables and clauses \textbf{MUST} match the
    header.

    \item All clauses, with the exception of the last clause, must end with
    ``0'', and can extend multiple lines.

    \item All variables are numbered $1$ through $N$, where $N$ is the number
    of variables. That is, if you say there are 10 variables, they must be
    labeled $1,2,\ldots,10$ and not $2,3,\ldots,11$ or any other convention.

    \item (For compatibility) Upon encountering a ``\%'' character outside of
    a comment, we will stop parsing immediately.
\end{enumerate}

Thus, valid input can take on many forms:

\begin{lstlisting}
p cnf 4 2
1 2 -3 0
-1 4 0
\end{lstlisting}

\begin{lstlisting}
4 2
1 2 -3 0
-1 4 0
\end{lstlisting}

\begin{lstlisting}
p cnf 4 2
c this comment doesn't have to be at the top
1 2 -3 0
-1 4 0
\end{lstlisting}

\begin{center}
\begin{lstlisting}
p cnf 4 2
c this comment doesn't have to be at the top
1 2 -3 0
-1 4
c the first two numbers are still the number
c of variables and the number of clauses
\end{lstlisting}
\end{center}

\begin{lstlisting}
p cnf 4 2 1 2 -3 0
-1 4 0
\end{lstlisting}

More information can be found here.\footnote{\url{https://people.sc.fsu.edu/~jburkardt/data/cnf/cnf.html}}

\section{Reading a CNF Equation Programmatically}

Now that we know what our input looks like, we can finally discuss reading it
programmatically. Firstly we must discuss how our SAT solver will be represented
in code. At the moment, we just need to keep track of

\begin{itemize}
    \item the number of variables
    \item the number of clauses
    \item each clause in our equation
    \item whether or not we've tried to solve the equation yet
    \item what our solution is
\end{itemize}

We also want some methods to parse the equation, construct a SAT solver,
solve the equation, print out the solution if it is satisfiable, and whether or not
there's a solution. Therefore, our object will look like so:

\begin{lstlisting}
public interface:
    SATSolver(string input)
    bool solveCNF()

private interface:
    void parse(string equation)
    int numVars
    int numClauses
    int[][] clauses
    optional<bool>[] vars
\end{lstlisting}

Our pseudocode is as follows:

\begin{lstlisting}
SATSolver(string input)
{
    vars = no-value

    if (input is file)
    {
        string content = readFile(input)
        parse(content)
    }
    else
    {
        parse(input)
    }
}
\end{lstlisting}

\begin{lstlisting}
parse(equation)
{
    // The method of parsing if there's a header
    if (contains(equation, "p cnf"))
    {
        int[1..N] Arr

        for each (line in input)
        {
            if (startsWith(line, 'c'))
            {
                skip to next line
            }
            else
            {
                for each (token in line)
                {
                    append(Arr, token)
                }
            }
        }

        numVars = Arr[1]
        numClauses = Arr[2]

        idx = 2

        for (i = 1 up to numClauses)
        {
            int[1..N] clause

            while (Arr[idx] != 0)
            {
                append(clause, Arr[idx])
                idx = idx + 1
            }

            append(clauses, clause)
        }
    }

    // Method of parsing when there's no header
    else
    {
        numVars = 0
        numClauses = 0

        int[1..N] Arr

        for each (line in input)
        {
            if (startsWith(line, 'c'))
            {
                skip to next line
            }
            else
            {
                for each (token in line)
                {
                    numVars = max(numVars, token)

                    if (token == 0)
                    {
                        numClauses = numClauses + 1
                    }

                    append(Arr, token)
                }

                idx = 0

                for (i = 1 up to numClauses)
                {
                    int[1..N] clause

                    while (Arr[idx] != 0)
                    {
                        append(clause, token)
                        idx = idx + 1
                    }

                    append(clauses, clause)
                }
            }
        }
    }
}
\end{lstlisting}

This code will be implemented slightly differently in the example code, as I
am using C++. There is also a big bug, which is that the code assumes that
because ``p cnf'' appears in the file, there is a header. This is not always true
because it does not consider whether or not the header is on the same line as a
comment, so this will break it:

\begin{lstlisting}
c p cnf
1 2 -3 0
-1 4 0
\end{lstlisting}

Because it assumes that there is a header when in fact there is not. I will
leave fixing this bug as an exercise to both the reader and the writer, however,
I will not be implementing this fix because if you're purposely trying to break
the constructor, you're not trying to solve a CNF equation anyways. The fix,
in case you're wondering, would be to just scan the lines and for each line that
doesn't start with ``c,'' check that it starts with ``p cnf,'' and if it does, then you
know for certain that there is a header. Now the real exercise is to determine
whether or not there's a header in just one pass instead of two.